\documentclass[../../main.tex]{subfiles}
\begin{document}
\section{微分形式的积分}\label{section:ch16-integration-of-differential-forms}

正如我们先在 $\mathbb{R}$ 中的区间上定义 $1$-形式的积分从而开始对线积分的讨论那样，我们在此先在 $\mathbb{R}^n$ 的适当子集上定义 $n$-形式的积分。眼下，让我们把注意力限于 $n\geqslant1$ 的情形。你应当确保自己熟悉多重积分的基本性质，这些性质总结于附录~C。

回忆一下，$\mathbb{R}^n$ 中的一个\emph{积分区域}是一个边界为零测集的有界子集。设 $D\subseteq\mathbb{R}^n$ 是一个积分区域，并设 $\omega$ 是 $\overline{D}$ 上的一个（连续的）$n$-形式。任何这样的形式都可以写成 $\omega=f\,\mathrm{d}x^1\wedge\cdots\wedge\mathrm{d}x^n$，其中 $f:\overline{D}\to\mathbb{R}$ 是某个连续函数。我们把 $\omega$ 在 $D$ 上的\emph{积分}定义为
\begin{align*}
\int_D\omega=\int_D f\,\mathrm{d}V.
\end{align*}
这可以更具提示性地写成
\begin{align*}
\int_D f\,\mathrm{d}x^1\wedge\cdots\wedge\mathrm{d}x^n=\int_D f\,\mathrm{d}x^1\cdots\mathrm{d}x^n.
\end{align*}
只要我们谈及在 $\mathbb{R}^n$ 的一个区域上积分一个 $n$-形式，就总是理解为这同一个积分，只是用形式记号表达。

这个定义的基本性质总结在下一个命题中。

\begin{proposition}[积分的线性]\label{proposition:ch16-linearity-of-integration}
设 $D\subseteq\mathbb{R}^n$ 是一个积分区域，$\omega,\eta$ 是 $\overline{D}$ 上连续的 $n$-形式，$a\in\mathbb{R}$。则
\begin{align*}
\int_D(\omega+\eta)&=\int_D\omega+\int_D\eta,\\
\int_D a\omega&=a\int_D\omega.
\end{align*}
\end{proposition}

\begin{proof}这可直接从多重积分对被积函数的线性性得到。
\end{proof}

\paragraph{积分区域}
我们主要关心的积分区域是那些具有分段光滑边界的区域。为了给出这个概念，设 $D\subseteq\mathbb{R}^n$ 是一个子集。我们称 $D$ 是\emph{恰当地位于其边界的一侧}，如果 $D$ 等于 $\overline{U}$，其中 $U$ 是某个开子集且 $\partial U=\partial D$。（例如，闭单位球 $\overline{\mathbb{B}}^n$ 就恰当地位于其边界的一侧，而 $\overline{\mathbb{B}}^n\setminus\mathbb{B}^{n-1}\subseteq\mathbb{R}^2$ 则不是。）

\begin{definition}[带分段光滑边界的区域]\index{区域!带分段光滑边界的}
一个紧子集 $D\subseteq\mathbb{R}^n$ 称为具有\emph{分段光滑边界}，如果它恰当地位于其边界的一侧，并且 $\partial D$ 是有限多个 $(n-1)$ 维光滑嵌入超曲面的并，这些超曲面（作为子集）在 $\mathbb{R}^n$ 中闭，并且两两之交是低维的。
\end{definition}

\begin{figure}[ht]
\centering
\includegraphics[width=0.6\textwidth]{fig-16-4.png}
\caption{具有光滑边界的有界子集}\label{figure:ch16-smooth-boundary}
\end{figure}

\begin{figure}[ht]
\centering
\includegraphics[width=0.55\textwidth]{fig-16-5.png}
\caption{边界不是分段光滑的有界子集}\label{figure:ch16-boundary-not-ps}
\end{figure}

\begin{figure}[ht]
\centering
\includegraphics[width=0.55\textwidth]{fig-16-6.png}
\caption{具有尖点的边界}\label{figure:ch16-cusp}
\end{figure}

这个定义的意图在于排除像\cref{figure:ch16-boundary-not-ps} 中那样的情形，那里边界不是分段光滑的；也排除像\cref{figure:ch16-cusp} 中那样的情形，那里边界含有尖点。另一方面，像\cref{figure:ch16-smooth-boundary} 中那样的区域确实有分段光滑边界。

具有分段光滑边界的区域是积分区域：这是\cref{proposition:ch06-graph-measure-zero} 及其推论的一个推论，因为一个具有分段光滑边界的紧子集的边界是有限多个光滑嵌入超曲面的并，而每一个这样的超曲面都有零测度。

接下来两个结果记录了积分在坐标变换和仿射子空间下的表现。它们的证明是初等的多重积分计算，在此略去。

\begin{proposition}[坐标变换下的积分]\label{proposition:ch16-change-of-coords}
设 $U,V\subseteq\mathbb{R}^n$ 是开子集，$G:U\to V$ 是一个保定向微分同胚。如果 $D\subseteq U$ 是一个积分区域，而 $\omega$ 是 $V$ 上一个紧支集 $n$-形式，则
\begin{align}
\int_{G(D)}\omega=\int_D G^*\omega.
\end{align}
若 $G$ 反定向，则等式右边带一个负号。
\end{proposition}

\begin{proposition}[仿射子空间上的积分]\label{proposition:ch16-affine-subspace-integration}
设 $S\subseteq\mathbb{R}^n$ 是一个 $k$ 维仿射子空间，带有从 $\mathbb{R}^n$ 继承的定向，$D\subseteq S$ 是一个相对于 $S$ 的积分区域。如果 $\omega$ 是 $S$ 上的一个连续 $k$-形式，则
\begin{align*}
\int_D\omega=\int_D f\,\mathrm{d}V_S,
\end{align*}
其中 $\omega=f\,\mathrm{d}V_S$，$\mathrm{d}V_S$ 是 $S$ 的标准欧氏体积形式。这个积分不依赖于 $S$ 的定向正交基的选择。
\end{proposition}

为了后面的目的，我们还需要把积分的定义稍稍推广到边界附近的情形。设 $D\subseteq\mathbb{R}^n$ 是一个具有分段光滑边界的有界闭区域。我们把 $D$ 的一个\emph{邻域}理解为 $\mathbb{R}^n$ 中包含 $D$ 的任意一个开集。如果 $\omega$ 是某个这样的邻域上的一个 $n$-形式，我们就定义
\begin{align*}
\int_D\omega=\int_D\omega|_D,
\end{align*}
只要该积分有定义；由上面的命题，这与用来定义 $\omega$ 的那个邻域无关。

\paragraph{有向流形上的积分}
现在设 $M$ 是一个 $n$ 维有向光滑流形（允许带边）。我们希望能够在 $M$ 上积分具有紧支集的 $n$-形式。关键事实是：$n$-形式在定向坐标卡之间的变换式自动地带有与定向相容的符号。

\begin{proposition}[形式积分的性质]\label{proposition:ch16-properties-of-form-integration}
设 $M$ 是一个 $n$ 维有向光滑流形（带边或不带边）。
\begin{enumerate}[(i)]
\item \textbf{线性：}如果 $\omega,\eta$ 是 $M$ 上具有紧支集的连续 $n$-形式，$a\in\mathbb{R}$，则
\begin{align*}
\int_M(\omega+\eta)=\int_M\omega+\int_M\eta,\qquad
\int_M a\omega=a\int_M\omega.
\end{align*}
\item \textbf{正性：}如果 $M$ 是定向的，且 $\omega$ 是一个非负的体积形式（即 $\omega_p$ 在每个保定向基上取正值），则 $\int_M\omega\geqslant0$，且等号当且仅当 $\omega\equiv0$。
\item \textbf{限制：}如果 $U\subseteq M$ 是一个开有向子流形，$\omega$ 是 $M$ 上一个支集含于 $U$ 的 $n$-形式，则
\begin{align*}
\int_M\omega=\int_U\omega.
\end{align*}
\item \textbf{微分同胚不变性：}如果 $F:M\to N$ 是一个保定向微分同胚，$\omega$ 是 $N$ 上具有紧支集的 $n$-形式，则
\begin{align*}
\int_N\omega=\int_M F^*\omega.
\end{align*}
若 $F$ 反定向，则右边带一个负号。
\item \textbf{可加性：}如果 $D_1,D_2\subseteq M$ 是具有分段光滑边界的区域，且 $D_1\cap D_2$ 不含 $n$ 维内部，则
\begin{align*}
\int_{D_1\cup D_2}\omega=\int_{D_1}\omega+\int_{D_2}\omega.
\end{align*}
\end{enumerate}
\end{proposition}

这些性质的证明都直接来自定义以及欧氏空间中积分的相应性质。

\begin{figure}[ht]
\centering
\includegraphics[width=0.6\textwidth]{fig-16-7.png}
\caption{流形上的链法则}\label{figure:ch16-chain-rule-manifolds}
\end{figure}

在实践中，积分常常是在由参数给出的区域上计算的。下面的命题把这一过程严格化。

\begin{proposition}[参数化下的积分]\label{proposition:ch16-integration-over-parametrizations}
设 $M$ 是一个 $n$ 维有向光滑流形（带边或不带边）。设 $U\subseteq\mathbb{R}^n$ 是一个开球或开立方体，$S\subseteq U$ 是一个具有分段光滑边界的有界闭区域，$F:U\to M$ 是一个光滑映射，在 $S$ 的邻域上单数，并且作为到其像的映射是保定向的。如果 $\omega$ 是 $M$ 上一个支集含于 $F(S)$ 的 $n$-形式，则
\begin{align*}
\int_M\omega=\int_S F^*\omega.
\end{align*}
若 $F$ 反定向，则右边带一个负号。
\end{proposition}

\begin{figure}[ht]
\centering
\includegraphics[width=0.6\textwidth]{fig-16-8.png}
\caption{用立方体参数化}\label{figure:ch16-cube-param}
\end{figure}

\begin{proof}因为 $F|_S$ 是到其像的一个嵌入，它有光滑的局部逆。这个结果于是可由\cref{proposition:ch16-change-of-coords} 以及单位分解得到。
\end{proof}

\begin{example}[曲面上 $2$-形式的积分]\label{example:ch16-surface-integral-2form}
设 $D\subseteq\mathbb{R}^2$ 是一个具有分段光滑边界的有界闭区域，$\sigma:D\to\mathbb{R}^3$ 是一个光滑嵌入，其像 $S=\sigma(D)$ 是 $\mathbb{R}^3$ 中的一个有向嵌入曲面。令 $\omega$ 是 $\mathbb{R}^3$ 上的 $2$-形式
\begin{align*}
\omega=x\,\mathrm{d}y\wedge\mathrm{d}z+y\,\mathrm{d}z\wedge\mathrm{d}x+z\,\mathrm{d}x\wedge\mathrm{d}y.
\end{align*}
则 $S$ 上 $\omega$ 的积分可以通过把它拉回到参数区域 $D$ 上来计算。写出 $\sigma(u,v)=(x(u,v),y(u,v),z(u,v))$，有
\begin{align*}
\int_S\omega=\int_D\sigma^*\omega,
\end{align*}
其中被积函数在标准笛卡尔坐标下展开后是一个普通的二重积分。
\end{example}

\begin{proposition}[紧 Lie 群上的 Haar 测度]\label{proposition:ch16-haar-measure}
设 $G$ 是一个紧 Lie 群。在 $G$ 上存在一个非零的、光滑的、左不变 $n$-形式 $\omega$（其中 $n=\dim G$），它在差一个常数因子的意义下唯一。用这样一个形式定义的积分 $\int_G\omega$ 称为 $G$ 上的一个\emph{左不变 Haar 积分}。它自动地也是右不变的。
\end{proposition}

\begin{proof}在李群的李代数 $\mathfrak{g}$ 上选取一个非零元素 $\omega_e\in\Lambda^n\mathfrak{g}^*$，并对每个 $g\in G$ 用左平移定义 $\omega_g=(L_{g^{-1}})^*\omega_e$。这就给出一个左不变的光滑 $n$-形式。它在差一个乘法常数的意义下唯一，因为 $\Lambda^n\mathfrak{g}^*$ 是一维的。紧性保证了 $\int_G\omega$ 有限。右不变性可由把右平移写成左平移与内自同构的复合、并注意到紧群上的内自同构在 $\Lambda^n\mathfrak{g}^*$ 上的作用是一个绝对值为 $1$ 的常数而得到。
\end{proof}

\begin{remark}
对于非紧的 Lie 群，类似的构造给出一个局部有限的左不变测度，但它一般不是右不变的。
\end{remark}
\end{document}
